\documentclass[11pt]{amsart}
\usepackage[T1]{fontenc}
\usepackage{lmodern}
\usepackage{amsmath,amssymb,amsthm,mathtools}
\usepackage{booktabs}
\usepackage{microtype}
\usepackage{xcolor}
\usepackage[colorlinks=true,linkcolor=blue!55!black,citecolor=blue!55!black,urlcolor=blue!55!black]{hyperref}

\newtheorem{theorem}{Theorem}[section]
\newtheorem{lemma}[theorem]{Lemma}
\newtheorem{proposition}[theorem]{Proposition}
\newtheorem{corollary}[theorem]{Corollary}
\theoremstyle{definition}
\newtheorem{definition}[theorem]{Definition}
\newtheorem{remark}[theorem]{Remark}

\newcommand{\Fix}{\operatorname{Fix}}
\newcommand{\Aut}{\operatorname{Aut}}

\title[Bounds for 17-point unlabeled posets]{Certified bounds for the number of unlabeled posets on 17 points}
\author{Zijian Zeng}
\address{Institute of Computer Science and Digital Innovation, UCSI University, Kuala Lumpur 56000, Malaysia}
\email{zijianzeng@foxmail.com}
\email{1002266693@ucsiuniversity.edu.my}
\date{August 31, 2026}
\hypersetup{pdftitle={Certified bounds for the number of unlabeled posets on 17 points},pdfauthor={Zijian Zeng}}

\begin{document}

\begin{abstract}
Let $U(n)$ denote the number of isomorphism classes of $n$-element
posets.  The exact value of $U(17)$ remains unknown.  We prove the explicit
partial result
\[
377293651326150106\leq U(17)\leq15754175376569459236.
\]
The lower bound is obtained from Burnside's lemma by constructing three
disjoint families of posets fixed by each permutation cycle type.  The upper
bound deletes a maximal point, then combines the complete 16-point
comparability-count histogram with the Kruskal--Katona clique bound.  We also
independently certify the known disconnected count
$4554367168841547$.  All arithmetic is integral, and independent evaluators
are supplied for both endpoints.  The interval is not an exact enumeration
of the 17-point classes.
\end{abstract}

\maketitle

\section{Introduction}

Write $U(n)$ for the number of finite posets on $n$ unlabeled points, where
only order-preserving isomorphisms are factored out.  Brinkmann and McKay
enumerated all classes through 16 points and obtained
\[
U(16)=4483130665195087
\]
together with extensive refined tables \cite{BrinkmannMcKay2002}.  This is
still the last term in the live table A000112 \cite{OEISA000112}.  A recent
preprint advances the labeled count to 19 by moment transfer, but explicitly
retains 16 as the unlabeled frontier \cite{Ayala2026}.  Labeled moments are
automorphism-weighted and do not determine the number of isomorphism classes.

Our contribution is a certified interval for the next unlabeled term.  It is
not a solution of the exact enumeration problem.  We also recertify the
published disconnected count of Mohammad, Noor, and Talukder
\cite{MohammadNoorTalukder2022}.  Searches of OEIS, arXiv, Crossref, zbMATH
Open, author pages, and exact result strings through August 31, 2026 did not
locate the displayed interval; this is a search-bounded statement, not a
claim that every possible source has been examined.

\section{Definitions and the disconnected part}

For a finite poset $Q$, let $d(Q)$ denote its number of order ideals,
equivalently its number of antichains.  Let $C(n)$ and $D(n)$ denote,
respectively, the numbers of connected and disconnected unlabeled
$n$-point posets.  Connectedness refers to the undirected Hasse graph (or,
equivalently, the comparability graph).  Unique decomposition into connected
components gives the Euler product
\begin{equation}\label{eq:euler}
1+\sum_{n\geq1}U(n)x^n=\prod_{k\geq1}(1-x^k)^{-C(k)}.
\end{equation}

\begin{proposition}\label{prop:disconnected}
The number of disconnected unlabeled posets on 17 points is
\[
D(17)=4554367168841547.
\]
More precisely, the numbers with $m$ connected components are as follows.
\end{proposition}

\begin{center}
\begin{tabular}{r@{\quad}r@{\qquad}r@{\quad}r}
\toprule
$m$ & count & $m$ & count\\
\midrule
2 & 4484639396830962 & 10 & 17905\\
3 & 68351780482060 & 11 & 2179\\
4 & 1340882409188 & 12 & 332\\
5 & 33946665077 & 13 & 64\\
6 & 1112290782 & 14 & 14\\
7 & 47328664 & 15 & 4\\
8 & 2624304 & 16 & 1\\
9 & 190010 & 17 & 1\\
\bottomrule
\end{tabular}
\end{center}

\begin{proof}
Adjoin a variable marking components in \eqref{eq:euler}.  The coefficient
of $x^{17}y^m$ in
\[
\prod_{k=1}^{16}(1-yx^k)^{-C(k)}
\]
counts the classes with exactly $m\geq2$ components.  Expanding the factors
with the certified values $C(1),\ldots,C(16)$ gives the displayed table, whose
sum is the asserted total.  Notice that $C(17)$ cannot occur in a term with
two or more components.  For independent checking, we also enumerated the
integer partitions $17=\sum k m_k$ and summed
\[
\prod_k \binom{C(k)+m_k-1}{m_k}.
\]
A third implementation first recovered $C(1),\ldots,C(16)$ from the published
$U(0),\ldots,U(16)$ by inverse Euler transformation.  All three computations
agree component by component.  The total matches the earlier published value
\cite[Table~4]{MohammadNoorTalukder2022}.
\end{proof}

\section{A constructive Burnside lower bound}

Let $P(r)$ be the number of partial orders on a fixed labeled $r$-element
set.  Fix $\sigma\in S_{17}$, and denote its distinguishable cycles by
$B_1,\ldots,B_r$, with lengths $\lambda_1,\ldots,\lambda_r$.

\begin{lemma}\label{lem:pairorbits}
The cyclic group generated by $\sigma$ has
$\gcd(\lambda_i,\lambda_j)$ orbits on $B_i\times B_j$.
Every nonempty such orbit projects onto both factors.
\end{lemma}

\begin{proof}
Simultaneous cyclic shift has orbit length
$\operatorname{lcm}(\lambda_i,\lambda_j)$.  Division of
$\lambda_i\lambda_j$ by this length gives the first assertion.  During a full
orbit the first and second coordinates each traverse their cycles.
\end{proof}

Uniformly replacing each point of a labeled $r$-point poset by the
corresponding antichain $B_i$ gives $P(r)$ distinct posets fixed by $\sigma$.
We next describe two further fixed families.

For a linear ordering $\pi$ of the blocks, choose between each adjacent pair
any nonempty union of its pair-orbits and orient it forward.  Put every point
of each earlier nonadjacent block below every point of the later block.  The
result is a poset: any two-step relation between adjacent blocks closes into
one of the complete nonadjacent relations.  Its block order and its adjacent
patterns are recoverable.  The number for a fixed $\pi$ is
\[
\prod_{i=1}^{r-1}
\left(2^{\gcd(\lambda_{\pi_i},\lambda_{\pi_{i+1}})}-1\right).
\]
The cases in which every adjacent relation is complete are indexed by all
permutations of the blocks; they are already uniform blow-ups.

For the second family, choose a nonempty matching $M$ on the cycle blocks.
Make each matched pair $i,j$ an isolated two-block component, choose either
orientation, and choose a nonempty proper union of its pair-orbits.  This
gives
\[
w_{ij}=2\left(2^{\gcd(\lambda_i,\lambda_j)}-2\right)
\]
possibilities.  Put an arbitrary labeled poset on the unmatched blocks and
blow it up uniformly.  The nonuniform block-pairs recover $M$, so the objects
are distinct.  When $r=2$ this family is omitted because it coincides with
part of the linearly ordered family; when $r>2$ its block quotient is
disconnected and hence cannot overlap that family.

For a partition $\lambda\vdash17$, let $m_j$ be the multiplicity of part $j$
and let $z_\lambda=\prod_j j^{m_j}(m_j!)$.  Define
\begin{align*}
T_\lambda={}&\sum_{\pi\in S_r}
\prod_{i=1}^{r-1}
\left(2^{\gcd(\lambda_{\pi_i},\lambda_{\pi_{i+1}})}-1\right)-r!,\\
M_\lambda={}&\sum_{\substack{M\text{ a nonempty matching}\\
                              \text{of the }r\text{ blocks}}}
P(r-2|M|)\prod_{\{i,j\}\in M}w_{ij},
\end{align*}
where $M_\lambda$ is set to zero for $r=2$.

\begin{proposition}\label{prop:fixedlower}
If $\sigma$ has cycle type $\lambda$, then
\[
\Fix(\sigma)\geq P(r)+T_\lambda+M_\lambda.
\]
\end{proposition}

\begin{proof}
The uniform, linearly ordered nonuniform, and matching constructions above
are fixed by $\sigma$.  They are internally injective.  A nonuniform relation
separates either additional family from the uniform family, while connected
versus disconnected block quotients separate the two additional families for
$r>2$.  The exceptional overlap at $r=2$ was removed by definition.
\end{proof}

\begin{proposition}\label{prop:lower}
One has
\[
U(17)\geq377293651326150106.
\]
\end{proposition}

\begin{proof}
The conjugacy class of type $\lambda$ has size $17!/z_\lambda$.  Exact
aggregation of Proposition~\ref{prop:fixedlower} over the partitions of 17
gives three contributions
\begin{align*}
L_0&=134198602505888509372423191552000,\\
L_1&=2210831977844367360,\\
L_2&=5971256590320944751628800.
\end{align*}
Thus the Burnside numerator is at least
\[
L=L_0+L_1+L_2=134198608477147310525345787548160.
\]
Burnside's lemma and integrality give
\[
U(17)\geq\left\lceil\frac{L}{17!}\right\rceil
=377293651326150106.
\]
\end{proof}

\section{A one-point-extension upper bound}

\begin{lemma}\label{lem:extension}
Let $\mathcal P_{16}$ be a set of representatives of the 16-point poset
classes.  Then
\[
U(17)\leq\sum_{Q\in\mathcal P_{16}}d(Q).
\]
\end{lemma}

\begin{proof}
Every 17-point poset has a maximal point.  Delete one and call the remaining
poset $Q$.  The strict lower set of the deleted point is an order ideal of
$Q$.  Conversely, adjoining a new maximal point above exactly an order ideal
defines a poset.  After transporting $Q$ to the chosen representative of its
class, this construction covers every 17-point isomorphism class.  Different
ideals can yield isomorphic extensions, which is why the conclusion is an
upper bound.
\end{proof}

Suppose $Q$ has $q$ strict comparable pairs.  Its antichains are the
independent sets of its comparability graph and hence the cliques, including
the empty set and singletons, of the complementary graph with
$N=\binom{16}{2}-q$ edges.  Write
\[
N=\binom{a}{2}+b,\qquad 0\leq b<a.
\]
The Kruskal--Katona theorem says that the colex graph maximizes the number of
cliques of every size; this graph formulation is recorded explicitly in
\cite[Proposition~1]{KirschRadcliffe2021}, based on the original theorem
\cite{Kruskal1963}.  The colex graph is a $K_a$, one further vertex joined to
$b$ clique vertices, and isolates.  Consequently
\begin{equation}\label{eq:imax}
I_{16}(q)=2^a+2^b+16-a-1.
\end{equation}
For $N=\binom{16}{2}$, the value is $I_{16}(0)=2^{16}$.

\begin{proposition}\label{prop:upper}
One has
\[
U(17)\leq15754175376569459236.
\]
\end{proposition}

\begin{proof}
Table~59 of Brinkmann--McKay \cite{BrinkmannMcKay2002} gives, for every
$0\leq q\leq120$, the exact number $h_q$ of 16-point unlabeled posets with
$q$ strict relations.  It checks that $\sum_q h_q=U(16)$.  Lemma
\ref{lem:extension}, the ideal--antichain bijection, and
\eqref{eq:imax} give
\[
U(17)\leq\sum_{q=0}^{120}h_q I_{16}(q)
=15754175376569459236.
\]
\end{proof}

Propositions \ref{prop:lower} and \ref{prop:upper} prove the interval in the
abstract.

\section{Exact verification and limitations}

The supplementary certificate consists only of integer input tables and
short Python programs using arbitrary-precision integer arithmetic.  The
disconnected count is evaluated by three methods.  The Burnside lower
numerator is evaluated both by a cycle-multiplicity recursion and by a
distinguishable-cycle bit-mask dynamic program.  The upper bound is evaluated
from the closed colex formula and, independently, by constructing all 121
colex graphs and testing all $2^{16}$ vertex subsets.  A separate parser
checks every serialized histogram entry against the archived published
table.  No floating-point value enters a theorem.

The exact value of $U(17)$ remains open.  In a complete Burnside computation,
permutations with several nontrivial cycles admit nonuniform relations not
exhausted by our constructed families.  Already cycle type $2^2 1^{13}$
introduces counts tied to marked cover relations over labeled 15-point
posets, a refinement not supplied by the labeled totals or antichain moments
used here.  Direct canonical generation also requires a nontrivial extension
of the public 16-point generator and a complete independently checked split
manifest.  Neither missing step is replaced by a finite test or runtime
estimate.

\section*{Data, code, and computational disclosure}

The supplementary archive contains the exact input certificates, all
independent verification programs, the archived-source comparison, and
documented one-command checks.  The 2026 labeled-moment residue table was also
reconstructed from its authors' public data as a frontier check, but it is not
used as a certificate for $U(17)$.  The mathematical claims were separated
from exploratory computations and checked by independent exact
implementations.  No proof assistant was used.

\bibliographystyle{amsplain}
\bibliography{references}

\end{document}
